\chapter{Introduzione}
\label{introduzione}

In questo capitolo si propongono degli esempi per gli oggetti utilizzati più di frequente in latex: la Sezione~\ref{citazioni} descrive come scrivere citazioni, la Sezione~\ref{oggetti-float} propone degli esempi di oggetti float, la Sezione~\ref{compilazione} descrive come compilare questo documento.

\section{Citazioni}
\label{citazioni}

Inserisco qualche citazione per mostrare la bibliografia. Per gli articoli accademici è quasi sempre possibile reperire i blocchi da inserire nel file bib da scholar, come ad esempio. Scholar in questo caso è una risorsa/sito online e per questo. Precediamo le citazione da uno spazio indivisibile tramite il carattere \textasciitilde.

\section{Oggetti float}
\label{oggetti-float}

Nella Sezione~\ref{figure-float} si propone un esempio di figura float, mentre nella Sezione~\ref{tabelle-float} si propone un esempio di tabella float.

\subsection{Figure}
\label{figure-float}

La Figura~\ref{fig:esempio} è un esempio di figura float.

\begin{figure}[htb]
    \centering
    \includegraphics[width=.4\columnwidth]{figures/example.pdf}
    \caption{Esempio di figura float in latex.}
\label{fig:esempio}
\end{figure}

\subsection{Tabelle}
\label{tabelle-float}

La Tabella~\ref{tab:esempio} è un esempio di tabella.

\begin{table}[htb]
    \centering
    \begin{tabular}{| c | l | r |}
        \hline
        allineamento centrale & allineamento a sinistra  & allineamento a destra
        \\
        \hline
        \hline
        centrale & sinistra & destra
        \\
        \hline
    \end{tabular}
    \caption{Esempio di tabella float in latex.}
\label{tab:esempio}
\end{table}

\section{Compilazione}
\label{compilazione}

Di seguito il codice da utilizzare per generare il pdf:
\begin{lstlisting}[language=bash]
$ pdflatex main.tex
$ bibtex main.aux
$ pdflatex main.tex
$ pdflatex main.tex
\end{lstlisting}

